\section{Related Work}
\label{sec:related}

\paragraph{Point-cloud video models.}
Static point-cloud networks~\cite{qi2017pointnet,qi2017pointnetpp,wang2019dgcnn,wu2024ptv3} were extended to point-cloud video by MeteorNet~\cite{liu2019meteornet}, point spatio-temporal convolutions (PSTNet~\cite{fan2021pstnet}), point 4D convolution with self-attention (P4Transformer~\cite{fan2021p4transformer}; PST-Transformer~\cite{fan2023psttransformer}), primitive-level modelling~\cite{wen2022pptr}, temporal point patches~\cite{benshabat2024_3dinaction} and the attention-free state-space Mamba4D~\cite{liu2025mamba4d}. None of them explains its decisions. The ``motion saliency'' of recent pre-training~\cite{wang2026mosaic} selects tokens to mask, and Mamba4D's ``disentangled'' and ``causal'' refer to efficiency and scan order, not to attribution.

\paragraph{Attribution for images, video and point clouds.}
Surrogates~\cite{ribeiro2016lime,lundberg2017shap} gave way to occlusion and mask optimisation~\cite{petsiuk2018rise,fong2019extremal}, gradients~\cite{selvaraju2017gradcam,sundararajan2017ig} and transformer-specific methods~\cite{abnar2020rollout,chefer2021transformer}; attention as explanation is contested~\cite{jain2019attention} and unavailable for state-space models. Video methods attribute frames~\cite{price2020playfair}, read spatio-temporal attention~\cite{wang2025staa}, learn temporal concept bottlenecks~\cite{knab2025motif,lee2025dance} or test concept sensitivity~\cite{videotcav2024}. The closest in spirit, DANCE~\cite{lee2025dance}, separates motion, object and scene \emph{concepts}, but it is ante hoc, works in a learned concept space and is correlational; \method intervenes post hoc on physical 3D regions of any trained model. All of these work on 2D RGB. Point-cloud saliency~\cite{zheng2019pointcloudsaliency}, gradient maps~\cite{tayyub2022ape}, occlusion maps for LiDAR detectors~\cite{schinagl2022occam} and interpretable classifiers~\cite{devries2025pointmil} explain a \emph{single} cloud with no time axis. Interpretability for dynamic scene \emph{representations}---language-embedded 4D Gaussians~\cite{fiebelman2025_4legs,li2025_4dlangsplat}, static/dynamic separation for editing~\cite{kwon2025instruct4dgs}, importance for pruning~\cite{zhang2024lp3dgs}, prototypes for static splats~\cite{galus2026xsplain}---serves grounding, editing or compression, not the attribution of a recognition decision.

\paragraph{Causal attribution and trustworthy evaluation.}
LENS~\cite{watson2021lens} operationalises Pearl's necessity and sufficiency for tabular and text data, CaCE~\cite{goyal2019cace} contrasts interventional effects with concept sensitivity, and LD-ViCE~\cite{ldvice2025} generates RGB video counterfactuals; ``motion attribution'' has been used for \emph{training-data} attribution in video generation (Motive~\cite{wu2026motive}). None of them reaches 3D geometry and motion, and none separates the two. Uncertainty methods for 4D occupancy~\cite{reliocc2025,kalble2025evocc,occuq2025} ask how reliable a prediction is, not what caused it. Whether an explanation can be trusted is itself an empirical question. Deletion-based faithfulness~\cite{petsiuk2018rise,deyoung2020eraser} tests whether highly ranked inputs matter to the model, and model randomisation~\cite{adebayo2018sanity} tests whether the explanation depends on what the model learned. Point-cloud XAI still lacks an agreed evaluation standard~\cite{chakraborty2026benchmark}. We therefore validate \method with all three kinds of evidence: faithfulness, sanity checks and ground truth.
